\chapter{从数字身体到真实身体：AgentOS 的跨身体工程验证路线}

\noindent\fbox{
\parbox{0.95\textwidth}{
\textbf{【本章总体说明】}\quad
本章讨论 AgentOS 从数字环境进入真实物理世界时的工程演化路线。这里我们不把“身体”狭义理解为机器人硬件，而将其定义为智能主体能够持续感知、预测、控制并接受现实反馈的代理承载界面，即

$$
\mathrm{Body}
=
\mathrm{ControlledAndSensedCarrierOfAgency}.
$$

因此，浏览器、桌面操作系统、Shell、移动操作机器人、陪伴机器人、车辆以及最终的人形机器人，并不是若干彼此孤立的产品形态，而是一组具有严格能力包含关系和验证依赖关系的 Body Scale 阶梯。我们将看到，从数字 Body 开始具有重要的方法论意义：它允许我们在低物理风险、强可观测性和高重复性的环境中首先验证 SGC、FCS、BPCC、KRA、技能沉降、现实回流以及 Kernel--Body Runtime 分工，再逐步增加空间连续性、动力学约束、身体内稳态、安全实时性、社会交互和全身自由度。由此，AgentOS 的具身路线不应被理解为“把一个越来越大的语言模型安装进越来越复杂的机器人”，而应理解为一个逐层扩大观测--行动闭环、逐层增加现实约束、逐层检验 Body Scale 的系统工程过程。其基本原则是：

$$
\boxed{
\text{先验证统一的主体动力学，再逐步扩大身体；先闭合低风险现实，再进入高代价现实。}
}
$$

}
}

\section{数字 Body 先行：以低风险环境验证统一智能主体}

从 AgentOS 的整体工程路线来看，数字 Body 并不是物理机器人出现以前的一种临时替代品，而是完整 Body 理论中的第一类合法身体。只要一个环境能够向 Agent 提供持续可观察状态，并允许 Agent 通过一组动作改变该环境，使动作结果重新成为下一时刻的观察，那么它就在功能意义上构成一个身体--世界闭环。对于某个 Agent $A$，设数字环境状态为 $x_t^{D}$，Agent 通过数字身体获得观察

$$
o_t^{D}=O_D(x_t^{D}),
$$

并输出动作

$$
u_t^{D}=\pi_A(h_t,E_t,\Theta_t,\Psi_t,o_t^{D}),
$$

环境随后按照

$$
x_{t+1}^{D}=F_D(x_t^{D},u_t^{D})
$$

发生改变。当新的 $x_{t+1}^{D}$ 再次通过观察接口进入 Agent 时，一个完整的

$$
\boxed{
World\rightarrow Body\rightarrow AgentKernel
\rightarrow Body\rightarrow World
}
$$

闭环便已经成立。这里世界可以是网页、操作系统、文件系统、应用程序、网络服务甚至一个大型软件工程环境，并不要求它必须由三维物理空间构成。

数字 Body 之所以应该位于 AgentOS 具身路线最前端，是因为它能够把“主体问题”和“物理机器人问题”暂时分离。现实机器人同时包含感知噪声、机械误差、通信延迟、执行器饱和、碰撞风险、电池状态、热约束和安全责任。如果在 AgentOS 尚未证明自己能够维持长期主体、正确调度工作记忆、形成技能、识别异常和处理现实反馈之前，就直接把全部问题叠加到机器人系统上，我们将很难判断失败究竟来自认知 Kernel、Body Runtime、传感器、控制器还是机械系统。因此数字身体实际上承担了一种实验上的变量隔离作用：我们首先把真实世界动力学中的大量不可控因素压低，使 AgentOS 可以在高可重复性的环境中验证自己的认知组织原则。

从 Body Scale 的观点来看，数字身体最重要的研究问题不是“Agent 是否会点击网页”，而是它能够验证多少未来身体都需要共享的结构。设从身体 $B_i$ 迁移到身体 $B_j$ 时，可直接复用的感知、影响、预测和控制结构分别为

$$
T_P(B_i,B_j),\qquad
T_F(B_i,B_j),\qquad
T_M(B_i,B_j),\qquad
T_C(B_i,B_j),
$$

则可以概念性定义身体迁移度

$$
\mathcal S_B(B_i,B_j)
=
\alpha_P T_P+
\alpha_F T_F+
\alpha_M T_M+
\alpha_C T_C.
$$

数字 Body 的战略意义，就是尽早验证哪些内部 CSO 和 Agent-CSO 可以保持稳定，而不依赖某一种具体身体。换言之，我们需要验证

$$
\boxed{
\frac{\partial \mathrm{Architecture}_{Agent}}
{\partial B_i}
\rightarrow 0
}
$$

这一架构目标：换身体应该主要改变身体接地与局部动力学，而不应该重新创造一个主体。

数字环境还具有另一个物理环境难以同时提供的优势，即几乎全部状态都可以被记录。Agent 的每一次观察、动作、工具调用、状态迁移、SPC 判断、TCE 调节、CCM 卸载以及技能调用都可以构成可重放轨迹

$$
\Gamma_D=
\{x_t,o_t,h_t,u_t,r_t\}_{t=0}^{T}.
$$

因此我们能够在相同环境快照上重复运行不同 Kernel 版本，也能够分析某次错误究竟来自错误观察、错误 Agent-CSO、错误规划还是控制执行。这使数字 Body 不只是工程验证平台，也成为 AgentOS 动力学的实验平台。

更重要的是，数字 Body 同样能够验证“现实权威”原则。即便一个网页或文件系统是数字世界，Agent 内部的预测也不能替代真实状态。Agent 认为文件已经保存，与文件系统实际确认写入成功是两回事；Agent 认为网页提交已经完成，与服务器真正返回成功状态也是两回事。因此：

$$
\boxed{
ActionIssued\neq ActionSucceeded
}
$$

在数字环境和物理环境中具有完全相同的意义。我们由此可以在进入机器人之前就建立 AgentOS 最关键的习惯：任何行动必须重新经过世界，任何世界状态必须通过 Body Runtime 回流，而不能由语言模型的内部叙事直接宣布成功。

因此，本书把数字 Body 放在真实机器人之前，并不是因为数字任务更容易，而是因为数字世界提供了一个可以精确拆解 AgentOS 理论假设的实验场。只有当 Persistent Agent 能够在数字身体中保持长期身份、稳定工作记忆、执行跨任务规划、形成可重用技能、识别自身错误并接受现实反证之后，我们才有理由增加身体动力学这一更昂贵的自由度。具身路线由此遵循的不是“模型能力最大化”，而是：

$$
\boxed{
\text{ClosedLoopCorrectness}
\rightarrow
\text{BodyComplexityExpansion}.
}
$$

\section{Browser Body：从网页工具调用到可持续的数字感知--行动闭环}

浏览器是 AgentOS 最适合作为第一种标准数字身体的环境之一，因为它同时具备视觉世界、结构化语义、持续状态、可逆动作和丰富现实反馈。传统 Computer Use 往往把浏览器理解成若干“点击、输入、滚动”的工具集合，而在 AgentOS 中，我们更希望把浏览器完整地理解为一种身体：网页是外部环境，Viewport 是局部感觉窗口，DOM、Accessibility Tree 和视觉图像是不同感觉通道，鼠标指针和键盘输入是数字效应器，下载文件、Cookie、Local Storage 与 Session 则属于身体和环境之间形成的持续状态。

设网页的真实状态为

$$
x_t^{web},
$$

它并不等同于当前屏幕像素 $I_t$，因为一个页面还具有不可见元素、滚动区域、网络加载状态、DOM 关系以及后端会话状态。因此 Browser Body 应形成一个多通道观察：

$$
O_t^{web}
=
\left(
I_t,\,
D_t,\,
A_t,\,
N_t,\,
H_t
\right),
$$

其中 $I_t$ 为视觉图像，$D_t$ 为 DOM 结构，$A_t$ 为 Accessibility 语义，$N_t$ 为网络/加载信息，$H_t$ 表示交互历史。Body Runtime 的任务不是把这五类原始输入全部交给 Kernel，而是把它们压缩为当前可行动现实：

$$
O_t^{web}
\xrightarrow{\mathrm{FBI}}
SGC_t^{web}\oplus FCS_t^{web}.
$$

在 Browser SGC 中，一个网页按钮不应只是一个二维 bounding box，而应是一个可行动实体：

$$
e_i=
(
id_i,\,
geometry_i,\,
semantic_i,\,
state_i,\,
affordance_i,\,
confidence_i
).
$$

例如一个“提交”按钮同时具有屏幕坐标、语义标签、当前 enabled/disabled 状态、所属表单、可能动作以及状态来源。这样，Kernel 操作的不再是“在 $(842,517)$ 点击一次”，而是：

$$
\mathrm{Activate}(Entity:\texttt{submit\_button}).
$$

前者把能力绑定在具体页面像素上，后者才形成可迁移的 Agent-CSO 与身体动作之间的接地关系。

Browser Body 因而是验证 SGC 的理想起点。它能够让我们测试一个核心问题：Agent 是否能够保持“对象恒常性”。网页滚动以后按钮位置发生变化，但对象本身没有改变；页面局部刷新后 DOM 节点编号改变，但语义对象可能仍然是同一个对象；页面跳转以后旧对象消失，新页面对象形成新的结构。因此 Browser Runtime 必须完成从 Frame-centric Representation 到 World-centric Representation 的转换。形式上可以写为：

$$
e_i^{t+1}
=
\operatorname{Track}
\left(
e_i^t,O_{t+1}^{web}
\right),
$$

而不是每一帧都重新把网页当成无历史的新图像。

浏览器也提供了验证 epistemic metadata 的良好环境。Agent 可以区分：

$$
Observed,\qquad
Derived,\qquad
Predicted,\qquad
Pending.
$$

例如“订单已经提交”只有在服务器响应或页面状态实际改变以后才能成为 Observed；点击之前，“提交后可能进入支付页面”只是 Predicted；点击完成但页面仍在加载时则是 Pending。通过这种区分，我们可以显著降低 Agent 中常见的“把计划当结果、把预测当事实”的错误。

与此同时，Browser FCS 可以开始验证现实对主体功能影响的表示。数字 Body 不会产生生物疼痛，但会存在自己的功能压力：页面加载超时、身份验证失败、剩余 token/计算预算下降、下载队列阻塞、频繁重试、权限不足、会话即将失效等。这些状态可以形成：

$$
FCS_t^{web}
=
[
urgency,\,
uncertainty,\,
resourcePressure,\,
controlInstability,\,
risk,\ldots
].
$$

它们不描述网页“是什么”，而描述当前网页交互正在怎样影响 Agent 的行动状态。于是 Browser Body 已经能够验证 SGC/FCS 的正交关系，而无须等待真实机器人。

浏览器还是技能编译的第一个重要试验场。第一次完成一个陌生网页任务时，大量结构可能进入工作记忆：

$$
h\rightarrow
\text{理解页面}\rightarrow
\text{规划步骤}\rightarrow
\text{逐步操作}.
$$

如果同类交互反复出现，稳定结构应逐渐迁移：

$$
AgentCSO_{explicit}
\rightarrow
AgentCSO_{procedural},
$$

最终形成 Browser Skill。此时 Kernel 不再逐点击规划，而是调用一个局部闭环：

$$
Skill_{login}
:
(SGC,FCS,Goal)
\rightarrow
ActionSequence,
$$

并仅在 residual 超过阈值时重新升级：

$$
\|\varepsilon_{web}\|>\theta
\Rightarrow
\mathrm{EscalateToKernel}.
$$

这正是“正常动力学留在局部，异常残差上浮”的第一种数字实现。

所以 Browser Body 的价值远远超过 Computer Use benchmark。它实际上是 AgentOS 第一个可以完整验证感知接地、语义世界模型、局部技能、现实验证、异常上浮以及长期主体连续性的身体环境。我们应当把浏览器看成一个微型具身世界，而不是一组 API。只有在这种意义上，Browser Agent 才真正开始接近 AgentOS 所说的“数字生命体”。

\section{Desktop Body：从单一网页扩展到持续存在的数字生活空间}

如果 Browser Body 验证的是一个受限应用空间中的具身闭环，那么 Desktop Body 所增加的关键变量就是\textbf{跨应用、跨窗口和跨资源域的持续世界结构}。在桌面环境中，Agent 不再只面对一个页面，而同时面对窗口、应用、系统通知、文件、剪贴板、输入设备、后台任务、外接设备和权限边界。因此 Desktop Body 是 AgentOS 从“工具操作器”向“数字空间持续主体”转变的重要一步。

设桌面世界状态为

$$
x_t^{desk}
=
(
W_t,\,
P_t,\,
F_t,\,
C_t,\,
D_t,\,
R_t
),
$$

其中 $W_t$ 表示窗口拓扑，$P_t$ 表示运行中的应用进程，$F_t$ 表示文件系统相关状态，$C_t$ 表示剪贴板和跨应用上下文，$D_t$ 表示设备和外设，$R_t$ 表示权限及系统资源状态。与 Browser 不同，Desktop 的重要特征在于：当前屏幕只展示整个数字世界的一小部分。因此 Agent 的可见界面

$$
I_t^{screen}
$$

只是对

$$
x_t^{desk}
$$

的局部投影。

这使 Desktop Body 首次逼迫我们认真处理“主体位置”。在物理世界中，人总是在某处；在数字世界中，Agent 的“位置”可以理解为当前焦点窗口、活动应用、工作目录、选中实体和可控制资源的组合：

$$
p_t^{digital}
=
(
window_t,
application_t,
workspace_t,
selection_t,
authority_t
).
$$

因此在 Desktop SGC 中必须显式表示：

$$
\boxed{
SelfInDigitalWorld
}
$$

而不只是 World。当前鼠标指针可以被视为最基础的 Digital End Effector，键盘焦点则是另一种动作通道；当前选区、拖拽对象和 active window 都属于身体与环境的接触关系。这也是为什么我们此前把 Cursor 视为数字身体末端执行器，而不是一个普通 UI 辅助符号。

Desktop Body 还会第一次显著暴露认知边界问题。一个真实桌面往往同时存在几十个窗口、成百上千个文件以及大量后台任务，如果全部状态持续进入 $h(t)$，工作记忆有限容量定理会立即被破坏。因此桌面环境要求 Body Runtime、Boundary、CCM 与 Scheduler 形成明确分工。粗略地可以写：

$$
x_t^{desk}
\xrightarrow{BodyRuntime}
\hat x_t^{desk}
\xrightarrow{Boundary}
x_t^{relevant}
\xrightarrow{CCM}
h_t.
$$

第一步把原始数字状态变成语义状态，第二步选择当前任务相关区域，第三步再根据工作记忆容量压缩进入显式认知面的结构。于是：

$$
\boxed{
DesktopScale\uparrow
\not\Rightarrow
hScale\uparrow
}
$$

反而越成熟的 Desktop Agent，越应该让绝大多数数字世界状态停留在低层或外部结构中。

Desktop Body 也是认知卸载理论最自然的实验场。文件系统并不只是被调用的工具，而是可以成为真正的外部认知结构。一个复杂计划可以从工作记忆卸载为文件：

$$
AgentCSO_{plan}^{h}
\rightarrow
ExternalCSO_{plan}^{file},
$$

未来通过索引重新激活：

$$
ExternalCSO_{plan}^{file}
\rightarrow
AgentCSO_{plan}^{h}.
$$

同样，表格、代码、数据库和绘图都可以视为内部结构向环境的沉降。于是 Desktop Agent 的成熟并不表现为“能够一直记住更多”，而表现为懂得什么时候把结构稳定地写进环境，并保持自己与这些外部结构之间的索引关系。

跨应用交互则进一步验证 Agent-CSO 与具体身体实现的分离。例如“把研究结果整理成报告并发送”这一目标并不属于 Word、Browser 或 Mail 中的某一个应用，它是一个更高层 Goal-CSO。Desktop Body 需要把它动态接地到多个身体功能：

$$
GoalCSO
\rightarrow
\{
Browser,
Editor,
FileSystem,
Mail
\}.
$$

如果未来换成另一套操作系统或另一组应用，高层 Goal-CSO 应保持基本不变，改变的是 KRA 与 Body Runtime 对这些高层结构的实现映射。这正是 Body Scale 在数字环境中的直接检验。

Desktop Body 还首次要求 Agent 形成真正的数字空间长期习惯。例如长期使用后，Agent 应逐渐知道某些项目文件位于什么位置、哪些应用适合完成哪些任务、用户的工作空间怎样组织、哪些目录不能修改，以及哪些操作需要确认。这些结构横跨 $E$、$\Theta$ 和外部文件系统，因此它们构成一个持续数字生活环境，而不只是一个“每次重新开始”的工具集合。

因此，从 Browser 到 Desktop 的跃迁，本质不是 UI 面积变大，而是世界结构从单场景变成多区域持续环境。Agent 必须开始处理位置、资源、边界、长期外部结构以及跨应用 Goal 的统一协调。只有当它能够在这种环境中长期保持行为连续，我们才真正获得了一个意义上较完整的 Digital Agent Body。

\section{Shell Body：以符号化高带宽动作验证程序性身体与外部认知结构}

Shell 是一种与浏览器和桌面截然不同的数字身体。Browser 和 Desktop 主要依靠空间化、视觉化和对象操作，而 Shell 把大量身体行为压缩成符号命令。一个复杂的文件移动、程序启动、数据处理甚至网络部署，可以通过一个短命令完成。因此 Shell Body 是研究“高层语义如何直接编译成高效程序动作”的理想环境，同时也是检验 AgentOS 认知卸载、技能编译和现实验证原则的高风险数字场景。

可以将 Shell 状态写为：

$$
x_t^{shell}
=
(
cwd_t,\,
env_t,\,
proc_t,\,
fs_t,\,
perm_t,\,
net_t
),
$$

其观测通常不是二维视觉图，而是结构化文本：

$$
o_t^{shell}
=
(
stdout_t,\,
stderr_t,\,
exitcode_t,\,
stateDelta_t
).
$$

动作则是命令：

$$
u_t^{shell}
=
Command_t.
$$

因此 Shell 的观测--行动闭环可以表示为：

$$
Command_t
\xrightarrow{OS}
StateTransition
\xrightarrow{Output}
Agent.
$$

它的一个重要性质是动作的“语义压缩率”极高。一次鼠标操作通常改变一个局部对象，而一条 Shell 命令可能改变成千上万个文件。因此 Shell Body 的动作空间具有高杠杆特征：

$$
\left\|
\frac{\partial x_{t+1}}
{\partial u_t}
\right\|
\gg
\text{ordinary GUI action}.
$$

这意味着 Agent 必须拥有更强的行动前预测和权限约束。

Shell 也使“计划”和“程序”的边界变得非常清楚。一个第一次执行的任务可能需要在 $h$ 中逐步推理，例如：

$$
\text{查找文件}
\rightarrow
\text{筛选}
\rightarrow
\text{转换}
\rightarrow
\text{移动}
\rightarrow
\text{验证}.
$$

当这一序列稳定以后，它完全可以沉降为一个 Script：

$$
AgentCSO_{explicit\,procedure}
\rightarrow
ExternalCSO_{script}.
$$

以后 Agent 只需要调用：

$$
u_t=Execute(script,\theta),
$$

其中 $\theta$ 是当前参数。也就是说，Shell 提供了 AgentOS 中“认知结构向可执行外部结构迁移”的最直接范例。程序不是与认知相反的东西，而是已经稳定到足以被编译和外化的认知结构。

然而 Shell 也特别容易暴露一个错误：Agent 可能在输出命令以后，立即在语言内部假设任务已经完成。正确的结构应该始终要求：

$$
Action
\rightarrow
OS
\rightarrow
Observation
\rightarrow
Verification.
$$

例如：

$$
\texttt{mv a b}
$$

仅仅表示 Agent 发出了修改文件系统的意图，并不证明目标已经成功。权限错误、文件锁定、磁盘故障、路径误解都可能导致不同结果。因此 Shell Body 是强化

$$
\boxed{
LearnedDynamics\neq RealityAuthority
}
$$

的重要环境。

由于 Shell 命令具有高杠杆作用，我们还必须在这里提前引入 Body Runtime 的权限包络思想。设动作 $u$ 作用于资源集合 $\mathcal R(u)$，其潜在影响为

$$
\mathcal I(u)
=
F(
|\mathcal R(u)|,
irreversibility,
privilege,
uncertainty
).
$$

当

$$
\mathcal I(u)>\theta_{authority},
$$

动作不能作为普通局部技能直接执行，而应进入更高层验证、用户授权或安全治理流程。由此可见，Shell 虽然没有物理碰撞，却已经具备真实意义上的高风险身体动作。它能够删除数据、修改权限、启动服务甚至影响其他机器，因此非常适合作为进入物理机器人前的“高后果数字行动实验场”。

Shell 同样能够强化外部世界与外部记忆的统一。文件系统中的脚本、配置、日志、模型、数据库和代码，本身都是 Universe-CSO 在数字环境中的结构。Agent 内部的 Agent-CSO 可以通过写代码重新组织这些外部结构：

$$
AgentCSO
\rightarrow
Code
\rightarrow
DigitalWorld'.
$$

而新的数字世界又通过输出和文件状态重新进入 Agent。于是：

$$
\boxed{
Mind\rightarrow Program\rightarrow World
\rightarrow Observation\rightarrow Mind
}
$$

已经在 Shell 环境中形成了和物理具身完全同构的结构闭环。

因此 Shell Body 的工程意义不是“让 LLM 学会 Bash”，而是验证一个更本质的问题：一个持续 Agent 是否能够把显式认知逐渐编译成高效程序结构，同时保持权限治理、现实验证、异常上浮和可恢复性。只有能够安全驾驭这种高带宽数字身体，AgentOS 才真正具备进入高自由度物理身体的认知基础。

\section{Mobile Manipulator：从数字状态空间进入连续物理空间}

Mobile Manipulator，即同时具备移动底盘和机械操作能力的机器人，是 AgentOS 从数字身体进入物理世界最合理的第一类综合平台。它比固定机械臂多出了空间导航和环境切换能力，又比完整人形机器人减少了双足平衡、全身协调和极端自由度，因此能够以较低系统复杂度验证几乎所有关键具身理论：身体中心 SGC、FCS、BPCC、Predictive Envelope、局部控制、Residual Escalation、技能编译以及现实回流。

进入物理世界以后，Body 状态可以写为：

$$
x_t^B
=
(
q_t,\dot q_t,
p_t^B,
v_t^B,
s_t,
e_t
),
$$

其中 $q_t,\dot q_t$ 表示执行关节状态，$p_t^B,v_t^B$ 表示身体姿态和速度，$s_t$ 表示内部资源及健康状态，$e_t$ 表示局部环境。真实动力学满足：

$$
\dot x_t^B
=
F_B(x_t^B,u_t,w_t),
$$

其中 $w_t$ 是无法完全建模的外部扰动。这里与数字世界最大的不同是：动作结果不再由确定的软件语义直接决定，而受到摩擦、惯性、接触、延迟、传感噪声和物体未知属性共同影响。因此 Kernel 无法以每一步语言推理直接控制执行器。

这正是 BPCC 存在的根本原因。设 Kernel 更新周期为 $\Delta$，BPCC 的局部更新周期为 $\delta$，则应满足：

$$
\delta\ll\Delta.
$$

Kernel 在较慢时间尺度上给出：

$$
\mathcal E_H
=
(
SGC^{target},
FCS^{constraint},
Goal,
ControlReference
),
$$

而 BPCC 在局部连续动力学中寻找：

$$
u_{t:t+\Delta}^{*}
=
\arg\min_{u}
\int_t^{t+\Delta}
L(x(\tau),u(\tau),\mathcal E_H)d\tau.
$$

高层并不告诉机械臂“每毫秒应该施加多少扭矩”，而是定义可行域、目标关系和约束；局部控制器则在这个可行域内部完成真正的身体闭环。因此：

$$
\boxed{
HigherLevelSetsFeasibleRegion,\qquad
LowerLevelOptimizesInsideIt.
}
$$

Mobile Manipulator 也是检验 SGC 从二维数字空间扩展到三维行动空间的关键平台。一个杯子不再只是视觉对象，它必须包含：

$$
Object=
(
Pose,
Shape,
Semantic,
Affordance,
Relation,
Confidence
).
$$

而机器人自身必须显式位于同一个 SGC：

$$
Body\in SGC.
$$

于是“桌上的杯子”真正意味着身体与杯子之间存在一个可达几何关系，而不仅是视觉标签。对于轮式机器人，楼梯可能是 Obstacle；对于有足机器人，同一楼梯可能是 Traversable。SGC 因而必须始终是：

$$
\boxed{
BodyCenteredActionableReality
}
$$

而不是抽象的客观世界数据库。

物理身体还第一次使 FCS 获得不可替代的地位。电池余量、电机温度、碰撞风险、关节负载、抓取稳定性、传感器可靠性等状态会直接影响当前可行动空间。例如：

$$
FCS_t
=
[
EnergyPressure,
ThermalPressure,
IntegrityPressure,
ControlInstability,
Threat,\ldots
].
$$

同一个任务在高电量与低电量状态下可能应产生不同策略，因此 Kernel 不能仅从 SGC 决定行动。真正的决策状态是：

$$
\boxed{
CognitiveInput
=
SGC\oplus FCS.
}
$$

Mobile Manipulator 还可以清晰验证异常上浮机制。正常导航时，BPCC 可以自行闭合：

$$
\|\varepsilon_B\|<\theta_1
\Rightarrow
NormalDynamicsStayLocal.
$$

如果局部模型反复失败：

$$
\|\varepsilon_B\|\geq\theta_2,
$$

则 Residual 被提升给 Kernel，由高层重新理解问题：道路可能被堵塞、物体可能比预期更重、抓取方式可能错误。这里需要强调，Residual 不只是一个数值 error，而是“当前局部功能已经无法解释现实”的信号。因此：

$$
\boxed{
ResidualEscalation
=
CrossScaleCognitivePromotion.
}
$$

同样重要的是能力沉降。第一次打开陌生门可能需要大量高层推理；经历多次成功轨迹以后，可以形成：

$$
ExplicitPlan
\rightarrow
TrajectoryMemory
\rightarrow
StructuralSkill
\rightarrow
BodySkill.
$$

最终 BPCC 或局部技能层能够执行大部分动作，高层只维持目标和异常监控。这种从显式控制到局部技能的迁移，是 AgentOS 能否真正解决长期具身计算成本问题的关键。

因此 Mobile Manipulator 不是“简化版人形机器人”，而是一个理论上极其纯净的具身验证平台：它已经拥有真正的三维世界、移动性、操作能力和物理风险，却仍然保持相对有限的身体自由度。它能够让我们在进入复杂人形身体之前证明：AgentOS 的高低层闭环是否真的能够在现实动力学中稳定耦合。

\section{Companion Robot：用长期共处检验身体、自我与生命周期的共同连续性}

如果 Mobile Manipulator 主要检验“一个 Agent 是否能够正确行动于物理世界”，那么 Companion Robot 检验的问题则更深：\textbf{同一个 Agent 能否长期生活在一个真实环境中，并让身体经验逐渐成为自己的生命周期结构。}这也是为什么陪伴机器人在 AgentOS 路线中具有远超过普通消费机器人的理论价值。它同时要求身体、记忆、自我、社会互动、内稳态、长期学习和身份连续，而这些能力无法通过一次性 benchmark 验证。

设 Agent 的长期状态为

$$
X_A(t)
=
[
h(t),E(t),\Theta(t),\Psi(t),\Omega(t),
X_B(t)
].
$$

对于普通短时机器人任务，我们往往只关心某个时间窗口

$$
t\in[t_0,t_0+\Delta T],
$$

但陪伴 Agent 的有效观察窗口可能扩展到：

$$
\Delta T
\sim
\text{months or years}.
$$

这意味着许多在短时测试中可以忽略的问题都会变成系统的主要变量。例如，一个临时错误是否被错误写入 E；重复互动是否改变 Θ；长期交往是否逐渐进入 Ψ；身体老化、维修和更换部件是否影响身份连续；某种错误自我解释是否形成长期偏置。

Companion Robot 因而是三相记忆真正进入现实的重要验证环境。每天的大量交互首先以工作状态进入 $h$，其中一部分形成：

$$
h\rightarrow E,
$$

反复出现的稳定结构进一步形成：

$$
E\rightarrow\Theta.
$$

例如某个家庭成员每天固定的作息、某个房间的空间布局、某种常见交互习惯，都不应该永远作为独立 Episode 重新处理，而应该逐渐沉降成结构先验。更慢地，哪些经历属于“我的经历”、哪些关系对主体长期重要，则会参与：

$$
E,\Theta\rightarrow\Psi.
$$

因此 Companion Agent 的价值就在于它迫使我们回答一个此前软件 Agent 很容易逃避的问题：记忆并不是无限积累日志，而是生命周期中的结构重组。

社会互动还会使 SGC 与 FCS 获得新的层次。一个家庭成员在 SGC 中不仅具有位置和身份，还具有长期关系：

$$
Relation(A,Human_i)=
[
history,
role,
trust,
interactionPattern
].
$$

与此同时，对方的出现可能改变 FCS 和 AHC，例如社会熟悉度、互动紧迫度、不确定压力等。这里我们仍然必须保持工程边界：这些变量代表功能性调制，而不直接等同于人类情绪体验。但是它们能够让同一个物理场景在不同社会历史下产生不同的行动倾向，这正是长期主体区别于无历史机器人控制器的关键特征。

Companion Robot 还首次使“身体历史”成为身份历史的一部分。设机器人更换某个执行器：

$$
B_t\rightarrow B_{t+1}.
$$

如果 Agent 的主体完全绑定在特定马达参数上，那么硬件维修就可能意味着新的 Agent；而 AgentOS 所追求的是：

$$
\Psi_{t+1}\approx\Psi_t
$$

在身体局部改变时保持高层身份连续，同时通过 KRA 与 Body Runtime 重新校准身体关系。因此陪伴机器人是 Body Scale 的长期检验：Agent 是否能够不断适应自己的身体变化，却仍然保持“同一个持续主体”。

此外，陪伴环境还天然提供了 Agent 培育的现实基础。一个刚部署的 Agent 不应立即拥有与成熟 Agent 相同的自治权限。随着：

$$
Experience
\rightarrow
Capability
\rightarrow
Evaluation
\rightarrow
Trust
$$

增加，其可行动范围才逐渐扩大：

$$
Authority_{t+1}
=
G(
Authority_t,
Capability_t,
Trust_t,
Risk_t
).
$$

因此 Companion Robot 同时把技术成熟、信任形成与自治成长连接在同一个生命周期上。

这也是为什么陪伴 Agent 比许多标准机器人 benchmark 更严格。一个机器人可以在实验室中以 $95\%$ 成功率完成抓取任务，但如果它在一年时间内不断积累错误记忆、身份漂移、社会关系混乱、控制偏置和恢复失败，那么它仍然不是成熟的 Persistent Agent。真正需要测量的不是单次成功率，而是：

$$
\boxed{
LongTermCoherence
+
Adaptation
+
Recoverability
+
IdentityContinuity.
}
$$

因此 Companion Robot 是 AgentOS 路线中的一个重要中间成熟点。到这里，身体不再只是执行任务的工具，而成为主体生命周期的一部分；环境不再只是场景，而成为长期生活空间；人与 Agent 的关系也不再只是输入输出，而成为能够进入记忆、自我和未来行为的持久结构。一个 Agent 能够长期“生活”以后，我们才有资格讨论更加高速、更危险和更受法规约束的身体。

\section{Vehicle Body：用高速、安全关键闭环检验多尺度控制的硬边界}

车辆身体与陪伴机器人代表了两种不同的复杂性。陪伴机器人强调长期社会性和生命周期，而车辆首先强调\textbf{高速、连续和安全关键的物理控制}。这使 Vehicle Body 成为检验 AgentOS 多尺度闭环架构是否具有真实工程严肃性的关键阶段。因为在这里，高层认知的任何延迟、错误自信或控制越权都可能快速转换为现实损失。

设车辆状态为：

$$
x_t^V=
[
p_t,v_t,a_t,\psi_t,
s_t^{vehicle},
e_t^{road}
].
$$

其动力学可表示为：

$$
\dot x_t^V
=
F_V(x_t^V,u_t,w_t),
$$

其中道路、天气、其他车辆和行人形成强时变扰动 $w_t$。车辆最直接揭示了时间尺度分离的必要性。底层姿态和车辆稳定控制可能运行在毫秒尺度；局部轨迹预测运行在几十到数百毫秒；路径规划和交通语义理解可能运行在秒级；长期路线选择、能源规划和用户目标又处在更慢层级。因此：

$$
\tau_{servo}
<
\tau_{BPCC}
<
\tau_{SGC}
<
\tau_h
<
\tau_{planning}.
$$

如果 Kernel 试图直接控制 steering torque、brake pressure 或轮端力矩，整个系统将跨越多个时间尺度而失去稳定性。

Vehicle Body 因而必须严格贯彻：

$$
\boxed{
CriticalSafety\nrightarrow CognitiveDependency.
}
$$

设 Safety Controller 为 $C_S$，BPCC 为 $C_B$，Kernel 为 $C_K$。正常情况下：

$$
C_K\rightarrow C_B\rightarrow Vehicle,
$$

但当安全边界

$$
h_s(x)\geq0
$$

即将被破坏时，$C_S$ 必须具有直接覆盖局部控制的权限：

$$
h_s(x_t)\rightarrow0
\Rightarrow
C_S>C_B>C_K.
$$

这里的“$>$”表示控制权优先级，而非功能价值大小。这体现了我们此前提出的：

$$
\boxed{
SafetyAuthority>BPCCAuthority.
}
$$

车辆也特别适合检验 Predictive Envelope。高层不直接输出一个精确控制序列，而输出目标走廊、速度区间、风险预算和禁止区域：

$$
\mathcal E_H^V
=
(
RouteCorridor,
VelocityEnvelope,
RiskBound,
ForbiddenSet
).
$$

BPCC 在其中进行短时轨迹优化：

$$
u^{*}
=
\arg\min_u
\sum_{k=0}^{N}
L(x_k,u_k)
$$

subject to

$$
x_k\in\mathcal E_H^V.
$$

这种结构使高层能够保持语义责任，例如“在下一个路口右转并保持安全距离”，同时低层根据实际道路不断重新调整连续动作。

Vehicle Body 对 SGC 的要求同样远高于普通目标检测。道路中的车辆、行人、交通信号、车道、遮挡区域和可行驶区域都必须存在于同一个预测性语义场中：

$$
SGC_t^V
=
Observed_t
\cup
Predicted_t
\cup
Uncertain_t.
$$

尤其重要的是 epistemic uncertainty。例如一个被大型车辆遮挡的区域，即便没有观测到行人，也不能被表示成“无人”，而应该是一块高不确定性空间。由此，SGC 从“我看见什么”进一步变成“当前我对现实知道什么以及不知道什么”。

FCS 则可以表达车辆自身对当前环境的功能状态，例如控制稳定裕量、轮胎抓地估计、碰撞风险、能源压力、传感器退化等：

$$
FCS_t^V
=
[
Threat,
ControlMargin,
SensorConfidence,
EnergyPressure,
Urgency
].
$$

这使 Kernel 能够理解“当前身体能不能完成某个计划”，而不是把 Body 当成无限执行能力的抽象 API。

车辆还是 Single Writer Principle 最严苛的场景之一。对于转向、制动等资源：

$$
\forall r_i,
\quad
|\operatorname{ActiveWriter}(r_i)|=1.
$$

Kernel、BPCC、安全控制器不能同时直接写执行器，否则会形成控制竞争。高层只能修改引用和约束，真正执行必须通过明确的控制权结构完成。

因此 Vehicle Body 的核心价值不在于商业场景本身，而是它逼迫 AgentOS 把此前大量认知理论接受控制工程的现实检验：时间尺度分离是否真实成立；预测是否与现实严格分离；安全是否真正独立于认知；异常是否能够在正确尺度升级；Kernel 是否知道自己的身体能力边界。如果这些问题无法在 Vehicle Body 中成立，那么所谓“通用具身智能”仍然只是语言层的想象。

\section{Humanoid Body：将人类环境的全部自由度纳入统一身体接口}

Humanoid 是 AgentOS Body Scale 路线中最复杂的综合身体，但它不应被理解为路线的起点。人形身体之所以具有特殊价值，并不只是因为它“像人”，而是因为人类文明的大量物理环境、工具、建筑、家具和工作流程本来就是围绕人类身体设计的。门把手高度、楼梯尺寸、桌椅结构、工具握柄和空间尺度都隐含了人类身体先验。因此，一个能够使用人类空间的 Humanoid 实际上获得了极其广泛的现有世界接口。

设人形机器人具有 $n$ 个可控自由度：

$$
q\in\mathbb R^n,
$$

完整身体状态为：

$$
x_t^H=
(q_t,\dot q_t,
c_t,
p_t^{base},
s_t^{int},
e_t),
$$

其中 $c_t$ 表示接触状态。与 Mobile Manipulator 相比，人形机器人的关键难度在于其动作不再局限于底盘移动加单臂操作，而是大量自由度通过接触动力学耦合。一个简单伸手动作可能同时涉及：

$$
\text{foot support}
+
\text{balance}
+
\text{torso}
+
\text{arm}
+
\text{grasp}.
$$

因此高层如果直接规划所有关节轨迹，将产生极高的认知和计算负担。

这恰好体现了 AgentOS 的成熟原则：

$$
\boxed{
MatureIntelligence=SelectiveExplicitCognition.
}
$$

对于成熟 Humanoid，大量身体动力学必须已经沉降到 Body Runtime。Kernel 不应该“思考如何保持平衡”，就像成年人走路时通常不会把每次踝关节修正提升到工作记忆一样。正常状态应满足：

$$
MostBodyDynamics\notin h.
$$

只有当局部控制不断产生无法解释的 residual，例如地面异常、身体损伤、陌生工具或动作失败时，相关 Agent-CSO 才被重新提升：

$$
LocalResidual
\rightarrow
h
\rightarrow
ExplicitReplanning.
$$

Humanoid 因而是技能沉降理论最严格的验证平台。我们可以把某项技能的成熟过程写为：

$$
h
\rightarrow
E
\rightarrow
\Theta
\rightarrow
BodySkill.
$$

例如第一次学习使用电钻，Kernel 可能需要显式处理姿势、按钮、工具方向和作用效果；随着经验增加，重复结构被压缩到 $\Theta$，最终形成 Body Skill，使低层能够快速执行稳定动作。高层只保留目标：

$$
Goal=\text{在指定位置钻孔},
$$

而局部身体系统处理力、姿态和接触。这就是 Agent-CSO 的功能迁移在真实高自由度身体上的完整实现。

Humanoid 还把 Body Scale 推到最强检验：一个同样的 Agent Kernel 是否能够从 Browser、Desktop、Mobile Manipulator 等身体中保留已经形成的高层世界结构。例如“门”“杯子”“工具”“危险”“等待”“帮助”这些 CSO 不应随着身体改变而全部重新学习。真正需要适配的是：

$$
CSO
\rightarrow
NewBodyGrounding.
$$

即：

$$
KnownMeaning
+
NewActionGeometry
\rightarrow
HumanoidCapability.
$$

如果所有知识在换成 Humanoid 后都必须重新从像素学习，那么系统实际上没有实现 Body Scale。

人形机器人还会放大 FCS 的重要性。随着身体结构复杂，局部温度、关节负载、碰撞、能量、平衡裕量和机械疲劳形成大量内部状态。若这些原始传感全部进入 Kernel，则会立即淹没 h。因此必须经过：

$$
RawBodyState
\rightarrow
FCS
\rightarrow
AHC/Kernel.
$$

例如几十个电机温度不必逐个进入认知层，而可以形成更高层的：

$$
ThermalPressure=0.72.
$$

这种压缩实际上使 FCS 成为物理身体内部复杂度向认知层的第一道结构屏障。

Humanoid 同样使社会环境成为身体环境的一部分。在家庭、工厂或公共空间中，一个动作不仅具有物理可行性，也具有社会意义。例如伸手接近一个人的身体在几何上可能可达，但不代表语义上允许。因此 SGC 必须同时包含：

$$
Geometry
+
SemanticRelation
+
SocialConstraint.
$$

这进一步证明，真正的具身世界模型不可能退化成三维占据地图。

从系统角度看，Humanoid 不是某一个模块的成功，而是所有先前理论共同工作的结果：

$$
\boxed{
Kernel
\bowtie
KRA
\bowtie
BodyRuntime
\bowtie
Reality.
}
$$

如果 Browser 验证语义身体，Desktop 验证持续数字空间，Shell 验证程序化高杠杆动作，Mobile Manipulator 验证三维具身闭环，Companion Robot 验证生命周期连续性，Vehicle 验证高速安全，那么 Humanoid 才是这些能力在一个高度通用身体上的最终汇合。因此，人形机器人应该被视为 Body Scale 的综合考试，而不是 AgentOS 的第一份作业。

\section{为什么不能直接从 LLM 跳到 Humanoid：跨尺度智能的工程依赖律}

在当前人工智能技术叙事中，一个非常具有诱惑力的想法是：只要基础模型足够强，再给模型接入摄像头、机械臂和机器人 API，就能够直接得到通用人形智能。然而从本书建立的智能动力学框架来看，这种路径忽略了一个根本事实：\textbf{语言模型的高层表征能力与持续具身主体所需要的多时间尺度闭环不是同一个工程对象。}从 LLM 直接跳到 Humanoid，并不是简单增加一个输出模态，而是试图一次跨越从符号推理到实时物理动力学之间的全部结构层级。

基础 LLM 的典型推理过程可以抽象为：

$$
z_{k+1}=F_{\theta}(z_k,input),
$$

其中 $k$ 更多表示推理步骤而非现实时间。物理机器人则必须满足：

$$
\dot x(t)=F(x(t),u(t),w(t)).
$$

前者允许暂停、重试和延迟，后者不会等待推理完成。机器人正在跌倒时，现实动力学依然继续：

$$
\frac{dx}{dt}\neq0
\qquad
\text{while LLM is thinking}.
$$

所以最基础的差别就是：

$$
\boxed{
InferenceTime\neq PhysicalTime.
}
$$

如果系统没有 BPCC、安全控制器和本地闭环，高层模型再强也不能取消这一事实。

第二个问题来自工作记忆有限性。Humanoid 同时存在视觉、听觉、触觉、关节、力矩、温度、电池、姿态、接触和社会环境等巨大状态空间：

$$
\dim(X_{raw})\gg\dim(h_{effective}).
$$

如果这些信号全部转成文本或 token 进入 LLM，系统实际上把本来应该在 Body Runtime 局部闭合的问题提升成全局认知问题。这不仅成本高，而且破坏了：

$$
\boxed{
NormalDynamicsStayLocal.
}
$$

成熟智能的方向恰恰相反：越稳定、越高频、越重复的结构越应该下沉，只有异常才占用稀缺工作记忆。

第三个问题来自现实权威。语言模型特别容易形成内部叙事连续性，而具身智能必须不断接受现实反证。Agent 认为手已经抓住杯子，与力传感器确认稳定抓取是不同事实；Agent 预测门会打开，与实际门轴被卡住也是不同事实。因此：

$$
\boxed{
Prediction
\neq
Observation,
\qquad
Action
\neq
Outcome.
}
$$

如果没有严格的 Reality Re-entry，机器人就可能依据自己的上一轮预测继续行动，最终使模型内部世界与物理世界逐渐分叉。

第四个问题来自身体迁移。如果我们直接针对 Humanoid 训练一个巨大端到端系统，那么很多世界知识与特定身体参数将被重新纠缠：

$$
WorldKnowledge
\otimes
BodySpecificPolicy.
$$

这使身体发生改变时难以确定哪些结构应该保留，哪些应该重新接地。AgentOS 的路线则希望形成：

$$
\boxed{
CSOOntology
\rightarrow
AgentCSO
\rightarrow
BodyGrounding.
}
$$

“杯子是什么”和“这只机械手如何抓杯子”应该属于不同尺度。前者可以跨身体复用，后者由 KRA、BPCC 和具体 Body Runtime 学习。只有这样，Body Scale 才真正存在。

第五个问题来自技能形成。一个没有局部编译机制的系统会出现极不自然的现象：第10000次走路仍然调用和第一次相同规模的语言推理。其平均认知成本近似：

$$
C_N
\approx
N C_{explicit}.
$$

如果能够将重复结构编译成技能：

$$
ExplicitAgentCSO
\rightarrow
CompiledBodySkill,
$$

则长期成本变为：

$$
C_N
\approx
N C_{local}
+
K C_{explicit},
\qquad
C_{local}\ll C_{explicit},
\qquad
K\ll N.
$$

真正成熟的 Agent 因此不是“什么都更会思考”，而是“大部分熟悉事情不再需要显式思考”。

第六个问题来自身份连续。直接把 LLM 当机器人“大脑”，往往会把一个 Session 的上下文误认为主体。实际上长期 Humanoid 必须跨越电源重启、任务切换、身体维修、软件升级和环境变化保持：

$$
\Psi_{t+\Delta}\sim\Psi_t,
$$

并让经历通过

$$
h\rightarrow E\rightarrow\Theta\rightarrow\Psi
$$

形成连续历史。否则系统拥有的是反复启动的高级控制器，而不是持续人工智能主体。

第七个问题来自安全与权限。Humanoid 的动作影响域远大于 Browser 或普通 Desktop。随着可控资源增加，行动潜在损失迅速增长。因此自治能力不能单纯由模型能力决定，而需要满足：

$$
Authority
=
F(Capability,Trust,Risk,Recoverability).
$$

低风险数字 Body 提供了一个逐渐建立这些控制机制的环境；如果直接进入 Humanoid，我们就同时把尚未充分验证的认知系统和高影响物理执行系统绑在一起。

基于以上原因，本章所提出的路线：

$$
\boxed{
DigitalBody
\rightarrow
Browser
\rightarrow
Desktop/Shell
\rightarrow
MobileManipulator
\rightarrow
CompanionRobot
\rightarrow
Vehicle
\rightarrow
Humanoid
}
$$

不应机械理解成每一个商业产品都必须严格按此顺序开发。它更重要的含义是一个\textbf{能力依赖偏序}。设身体平台集合为 $\mathcal B$，若平台 $B_j$ 的稳定实现依赖平台 $B_i$ 已验证的一组能力，则定义：

$$
B_i\prec B_j.
$$

真实工程路线实际上是在这张偏序图上寻找满足成本、风险和验证效率约束的路径：

$$
\Gamma^*
=
\arg\min_{\Gamma}
\left[
Cost(\Gamma)
+\lambda_R Risk(\Gamma)
+\lambda_U Uncertainty(\Gamma)
\right],
$$

subject to

$$
Dependencies(\Gamma)=Satisfied.
$$

因此“数字身体先行”并不是商业教条，而是一种最小化未验证变量同时进入系统的方法。

当我们沿着这条路线逐步前进时，真正不断扩大的是 Agent 与现实的耦合空间：

$$
\mathcal A_{Browser}
\subset
\mathcal A_{Desktop}
\subset
\mathcal A_{Physical}
\subset
\mathcal A_{Humanoid},
$$

同时现实约束也越来越强：

$$
Risk_{Browser}
<
Risk_{Manipulator}
<
Risk_{Vehicle/Humanoid}.
$$

每扩大一层 Body Scale，我们都要求此前已经形成的高层 Agent-CSO、身份、自我、记忆和认知控制尽可能保持连续，只把新的身体动力学留给新的 Body Runtime 去学习。于是身体扩张不再意味着重新训练一个 Agent，而意味着：

$$
\boxed{
SamePersistentAgent
+
NewEmbodiedGrounding.
}
$$

从这个意义上，本章所描述的路线最终指向的是一种与传统“机器人模型升级”完全不同的发展方式。我们不是一次性制造一副最复杂的身体，再希望一个大模型能够适应它；而是在一个持续主体保持身份连续的条件下，让它从数字世界开始，逐步获得更丰富的感知、更大的行动空间、更快的控制回路、更高的现实风险以及更复杂的社会关系。随着身体不断扩展，过去形成的认知结构不应被清零，而应成为下一种身体学习的先验。

因此，AgentOS 的具身路线最终可以浓缩为：

$$
\boxed{
\textbf{
不是让一个模型突然进入人形身体，
而是让一个持续智能主体逐步扩大自己能够感知、预测、控制和承担后果的现实边界。
}
}
$$

而这正是 Body Scale 在整个 AgentOS 工程路线中的真正意义：身体不是智能的外围设备，身体扩张本身就是智能体生命周期中对现实可达空间的一次次扩张；每一次扩张，都必须以此前已经稳定形成的主体结构、认知结构和控制闭环为基础，而不能用更大的模型规模替代这些缺失的动力学层级。
